\section{The Anti-Substitution Principle}
\label{sec:anti-substitution}

The responsibility architecture governs harms produced through AI deployment. A related problem arises when deployment becomes an asserted reason to withdraw responsibility altogether. A public agency points to an automated portal instead of providing the process its benefit scheme requires. A care institution points to a companion bot instead of staffing required support. A professional points to a model output instead of exercising the judgment attached to the professional role. In each setting, the technology may perform valuable work. The legal error is to infer from availability that an existing duty has been satisfied.

The \term{anti-substitution principle} prevents that inference:

\begin{quote}
Where law independently imposes a duty of care, process, accommodation, professional judgment, or public service, the availability or use of AI does not by itself discharge or narrow that duty. Automation may perform part of the duty when the governing law permits and the deployment preserves the protection the duty exists to secure.
\end{quote}

This is a rule of performance, not a preference for manual administration. An \term{automatic instrumentality} can perform every operational step needed to satisfy a duty; its availability alone is not performance. The court asks whether the legally protected function remains available through the deployed system and whether an answerable institution maintains correction when the automated path fails. That formulation permits AI to improve access, consistency, or capacity while keeping the duty attached to the person or organization law already made answerable.

\subsection{A Rule Courts Can Apply}

Four questions turn the principle into a decision rule.

First, \emph{what is the independent source and scope of the duty?} The source may be statute, regulation, contract, custody, guardianship, professional law, a tort special relationship, or constitutional process attached to a protected interest. The principle does not manufacture that predicate. It prevents a duty holder from treating automation as an implicit amendment to it. Constitutional due process, for example, supplies no general affirmative obligation to provide care; where legislation creates an enforceable entitlement, the text and remedial structure define the protection.\lrfootnote{See \casecite{deshaney1989}; cf. \casecite{talevski2023} (recognizing section 1983 enforcement of two specified Federal Nursing Home Reform Act rights).}

Second, \emph{which performance does the duty require?} Law often protects a function rather than a particular staffing form. Notice must convey the material basis for a decision. Review must reach an actor with information and authority to correct error. Professional judgment must integrate the considerations the governing standard makes relevant. An accommodation must afford meaningful access. Care obligations may require monitoring, response, continuity, or safe transition. Stating the protected function keeps the inquiry from collapsing into either ``a human must do every step'' or ``a chatbot exists, so the duty is complete.''

Third, \emph{does the deployed process preserve that function in practice?} Relevant evidence includes whether the system can handle foreseeable exceptions; whether an affected person receives intelligible notice; whether accessibility and language needs are met; whether the person can contest the material basis; whether uncertainty and limitations reach a competent reviewer; whether the reviewer has time, authority, and an effective intervention; whether failure triggers escalation; and whether the institution monitors outcomes after release. These are cascade variables. Functional equivalence is demonstrated through the control path, not through the label ``human in the loop.''

Fourth, \emph{who remains responsible for performance and correction?} Delegating a task to software does not relocate the duty to the software. The institution must name a controller who can restore service, decide an exception, correct a record, or arrange an alternative when the automated path fails. The responsibility charter supplies that allocation, while the production rule preserves the evidence needed to test it. Where the deployed process performs the legally required function, the duty is satisfied through automation. Where it does not, technological novelty is neither performance nor excuse.

The rule gives a court a narrow holding. The court need not decide that automated service is categorically inferior, that every claimant has a right to personal attention, or that one model design is mandatory. It can identify the predicate duty, specify its protected function, examine the deployed control path, and determine whether the responsible institution provided the performance law requires.

\subsection{Public Administration and the Right to an Effective Path}

Public-benefits cases show the rule already latent in doctrine. In \emph{K.W. ex rel. D.W. v. Armstrong}, Idaho used a budget tool to allocate community-based services to people with developmental disabilities. The notices did not adequately disclose how individualized budgets were calculated or permit recipients to identify and challenge errors. The Ninth Circuit affirmed preliminary relief on procedural-due-process grounds.\lrfootnote{\casecite{kWArmstrong2015}.} In \emph{Arkansas Department of Human Services v. Ledgerwood}, changes associated with an assessment methodology reduced in-home attendant-care hours, and the Arkansas Supreme Court upheld temporary injunctive relief tied to the agency's failure to follow required rulemaking procedure.\lrfootnote{\casecite{ledgerwood2017}.}

Neither decision forbade algorithmic allocation. Each required automation to remain inside the law governing the benefit. The agency could use a model to improve consistency or direct scarce resources, but the model could not replace the notice, explanation, rulemaking, or contest path the legal scheme required. Danielle Citron's account of technological due process explains the general mechanism: automated public systems can encode policy and error at scale while notice and hearing remain designed for discrete human decisions.\lrfootnote{See \artcite{citronTechDueProcess2008}.} The operative response is to preserve an intelligible basis, an opportunity to contest, and authority to correct before a system-level error becomes a durable administrative rule. Those protections inherit their force from existing public law, including the due-process principles exemplified by \emph{Goldberg v. Kelly}.\lrfootnote{See \casecite{goldberg1970}.}

The anti-substitution principle adds a control test to that doctrine. A portal that explains a decision and routes a supported challenge to an empowered official may improve process. A chatbot that repeats the result, cannot expose the material basis, and loops the beneficiary back to itself does not create review merely because it responds in natural language. The difference lies in the upward feedback and correction authority. A public body remains the duty holder and must be able to show where the path terminates.

The same analysis applies to capability tiering. A government may restrict tools, data, or model access for security, competence, or resource reasons. If a lower tier determines access to an existing entitlement, the affected person must retain whatever explanation, review, and accommodation the governing law requires. The availability of a general public chatbot cannot substitute for an appeal channel, and a hidden automated restriction cannot substitute for the accountable decision the statute assigns to an official.

\subsection{Care, Professional Judgment, and Relational Systems}

Medical AI illustrates how the principle supports high-value automation. A decision aid may outperform a clinician on a bounded prediction, extend expertise to an underserved setting, or detect a pattern a person would miss. The professional duty can nevertheless require facts outside the model, patient preferences, uncertainty, informed consent, follow-up, and intervention when the workflow fails. A 2026 medical-AI consensus accordingly combines access evaluation and real-world validation with multidisciplinary review, traceability, continuing monitoring, competence, and circuit breakers.\lrfootnote{See \artcite{caoMedicalConsensus2026}.} Recent scholarship identifies four conditions for meaningful oversight---epistemic capacity, cognitive space, decisional authority, and intervention effectiveness---and treats oversight as a property of the clinical system rather than a ceremonial human presence.\lrfootnote{See \artcite{vanDeSandeOversight2026}; see also \bluecite{whoEthicsAIHealth2021}.}

Under the four-question rule, the institution first states the actual duty: for example, competent professional judgment and follow-up under the applicable law. It then identifies which functions the model performs and which remain elsewhere in the workflow. The cascade shows whether the clinician receives enough information, time, and authority to act; the charter shows who responds when performance drifts or an exception exceeds the model's scope. AI can carry a large share of the work. Responsibility stays with the professional and institution that law authorizes and regulates.

Relational AI presents a different opportunity and the same legal structure. A companion can reduce isolation, translate, remind, help a user rehearse a difficult exchange, or extend supportive contact. Its availability does not authorize a care institution to withdraw staffing required by another law, an insurer to deny covered treatment, or a public program to replace an accessible service path with an endless automated loop. The relevant comparison is not human versus machine in the abstract. It is the protection promised by the governing duty versus the protection actually delivered through the deployment.

The principle also disciplines provider-controlled transitions. A voluntary companion service does not become property owned by the user, and a provider does not incur a general obligation to operate every persona indefinitely. But where consumer, disability, health, contract, or another law independently recognizes a relevant duty, foreseeable dependence can make continuity notice, data export, crisis safeguards, and a safe exit part of its performance. Research on companion loss and repeated relational exposure makes those consequences predictable enough to enter the process model.\lrfootnote{See \bluecite{defreitasReplika2025}; \artcite{banksDeletion2024}; \bluecite{kirkSteering2026}.} The protected interests are the user's autonomy, psychological health, safety, and access to an intelligible transition---not ownership of an artificial relationship.

\subsection{Automation That Performs the Duty Is Not Substitution}

The principle's boundary is affirmative. It protects automation that actually performs the legal function. A multilingual benefits assistant that makes notices intelligible, a clinical system that gives a professional better evidence and usable override authority, or a companion service that supplements human support can strengthen the underlying right or duty. Requiring the institution to name the function, preserve feedback, and maintain correction makes those deployments more dependable.

Two distinctions keep the rule exact. The first separates a protected function from a demanded medium. Unless the governing law says otherwise, there is no general right to a human conversation, a particular employee, a favorite model, or permanent access to a commercial persona. The second separates human responsibility from human tokenism. Preserving a duty does not require a person to repeat the model's work. It requires an institutional path by which the beneficiary can understand, challenge, and obtain the performance the duty secures.

The anti-substitution principle therefore completes the responsibility architecture. The front-door rule prevents an enterprise from pointing inward to an unidentified worker or outward to an artificial agent. Anti-substitution prevents a duty holder from pointing to the mere existence of an automated service. In both settings, the cascade asks the same question: where is the actor with the information and authority to make the legal protection effective?
